\documentclass[10pt,a4paper]{amsart}
\usepackage[margin=19mm]{geometry}
\usepackage{amsmath,amssymb,mathtools}
\usepackage[scr=rsfs,cal=euler]{mathalfa}
\usepackage{xfrac}
\usepackage[unicode,hidelinks,bookmarks=false]{hyperref}
\newcommand{\ie}{i.e.\ }
\newcommand{\cf}{cf.\ }
\newcommand{\resp}{respectively\ }
\newcommand{\Subsection}{\subsection}
\providecommand{\nobreakdash}{\nobreak\hbox{-}}
\setlength{\emergencystretch}{2em}
\multlinegap0pt
\usepackage{amssymb}
\usepackage{mathtools}
\usepackage[shortlabels]{enumitem}
\usepackage{tikz}
\newtheorem{theorem}{Theorem}
\DeclareMathOperator{\Exp}{Exp}
\newcommand{\PP}{\mathbb{P}}
\title{The Dual-Population Benchmark Model}
\author{Alexander Gnedin}
\date{Source: arXiv:2608.15769v1 (CC BY 4.0)}
\begin{document}
\maketitle

\noindent\emph{Source and licence.} The Dual-Population Benchmark Model. Version arXiv:2608.15769v1 (CC BY 4.0), distributed by arXiv under the Creative Commons Attribution 4.0 International licence, http://creativecommons.org/licenses/by/4.0/, as recorded in the arXiv licence metadata for this exact version and shown on the abstract page https://arxiv.org/abs/2608.15769v1. Full licence notice: \texttt{licenses/arxiv-2608.15769-CC-BY-4.0.txt}.\par\smallskip

\noindent\textit{Editorial scope.} Counted unit: the article's theorem prop:coalescent-blocks (source0.tex lines 733--736). The article's complete original proof of it is delivered with it (source0.tex lines 738--773, one \texttt{proof} environment of 36 lines). No mathematics is written, repaired or re-derived: the statement and the proof are the article's own lines, in the article's own notation, and no retained line is substituted or re-flowed.  No further source-local context is needed: every cross-reference of every retained line resolves inside the two delivered spans.  Nothing was omitted from any retained span, so the package carries no omission record.  No retained line contains a citation, so the wrapper carries no bibliography and no bibliography entry.  No retained line contains a figure environment, an image-inclusion command or drawing code, so no image is delivered and the package declares no asset dependency.  Editorial formatting only, in addition: an AMS \texttt{amsart} preamble that restates the article's own declarations and macros used by the retained lines, namely the environments \texttt{theorem} on separate counters, together with the standard packages those lines need.  The retained environments print in this document as Theorem 1 in order of appearance, whereas the article prints the same items with the numbering of its own full document.  The complete native source of the article as distributed in the arXiv e-print for this exact version is bundled unchanged as \texttt{sources/207782/source0.tex}.
\begin{theorem}[Coalescent block sizes via matched-rate benchmarks]
\label{prop:coalescent-blocks}
Let \(B_1,\dots,B_k\) be i.i.d.\ \(\Exp(1)\) benchmarks and let the future sample of size \(n\) be generated by \({\rm S}\). Conditionally on every one of the \(k+1\) cells being nonempty, the sizes \((A_0,\dots,A_k)\), regarded as an unordered multiset attached to specific future individuals, have  the same distribution as the block sizes of Kingman's \(n\)-coalescent at the moment it has \(b=k+1\) blocks.
\end{theorem}

\begin{proof}
Write \(b=k+1\). It suffices to compare, for every specific labelled partition of \(\{1,\ldots,n\}\) into blocks of sizes \(n_1,\ldots,n_b>0\) (summing to \(n\)), the probability each model assigns to it.

\emph{Benchmark model.} Conditioned on nonemptiness, \((A_0,\ldots,A_k)\) is uniform over the \(\binom{n-1}{b-1}\)  compositions of \(n\) with \(b\) positive parts, and given the composition, the assignment of specific future individuals to specific cells is uniform over the \(n!/\prod n_i!\) possibilities. A specific labelled \emph{set} partition corresponds to \(b!\) such outcomes (one for each assignment of the \(b\) specific blocks to the \(b\) ordered cells), each of probability \(\bigl[1/\binom{n-1}{b-1}\bigr]\cdot\bigl[\prod n_i!/n!\bigr]\). Hence
\begin{equation}
\label{eq:cbench}
\PP_{\rm bench}(\text{specific partition})=c_b^{\rm bench}\prod_i n_i!,\qquad
c_b^{\rm bench}=\frac{b!}{n!\binom{n-1}{b-1}}.
\end{equation}

\emph{Coalescent.} We show by downward induction on \(b\), from \(b=n\) to \(b=1\), that 
$$\PP_{\rm coal}(\text{specific partition})=c_b^{\rm coal}\prod_i n_i!$$
for a constant \(c_b^{\rm coal}\) depending only on \((n,b)\), not on the partition. At \(b=n\), the unique partition into singletons has probability \(1=c_n^{\rm coal}\), so \(c_n^{\rm coal}=1\). For the inductive step, fix a target partition at \(b-1\) blocks of sizes \(n_1,\ldots,n_{b-1}\). Since the coalescent merges a uniformly random \emph{pair} of the \(b\) current blocks, independently of block sizes, the target is reached exactly when, at the previous step, some target block \(i\) was instead present as two separate parts \((A,B)\) with \(A\sqcup B\) equal to that block, and this specific pair merged. An elementary count gives, summed over the unordered nonempty splits of an \(m\)-set,
\[
\sum_{A\subsetneq[m],\,A\ne\emptyset}\tfrac12\,|A|!\,(m-|A|)!=\tfrac12\sum_{j=1}^{m-1}\binom mj j!(m-j)!=\tfrac12(m-1)\,m!.
\]
Summing over which of the \(b-1\) target blocks was the freshly merged one,
\[
\PP_{\rm coal}(\text{target})=\frac1{\binom b2}\sum_{i=1}^{b-1}c_b^{\rm coal}\Bigl(\prod_{j\ne i}n_j!\Bigr)\cdot\tfrac12(n_i-1)\,n_i!
=c_b^{\rm coal}\Bigl(\prod_jn_j!\Bigr)\cdot\frac{\sum_i(n_i-1)}{2\binom b2},
\]
and since \(\sum_{i=1}^{b-1}(n_i-1)=n-(b-1)\), this gives a recursion independent of the target partition:
\begin{equation}
\label{eq:crecur}
c_{b-1}^{\rm coal}=c_b^{\rm coal}\cdot\frac{n-b+1}{b(b-1)},\qquad c_n^{\rm coal}=1.
\end{equation}

\emph{Matching.} The constant \(c_b^{\rm bench}\) of \eqref{eq:cbench} satisfies \(c_n^{\rm bench}=1\) (immediate at \(b=n\)) and, using  \(\binom{n-1}{b-1}=\binom{n-1}{b-2}\cdot\frac{n-b+1}{b-1}\),
\[
c_b^{\rm bench}\cdot\frac{n-b+1}{b(b-1)}
=\frac{(b-2)!}{n!}\cdot\frac{n-b+1}{\binom{n-1}{b-1}}
=\frac{(b-2)!}{n!}\cdot\frac{b-1}{\binom{n-1}{b-2}}
=\frac{(b-1)!}{n!\binom{n-1}{b-2}}=c_{b-1}^{\rm bench},
\]
giving recursion~\eqref{eq:crecur}. Since \(c^{\rm coal}\) and \(c^{\rm bench}\) satisfy the same recursion from the same boundary value at \(b=n\), they coincide for every \(1\le b\le n\); the two laws on specific labelled partitions -- hence on unordered block sizes -- are identical.
\end{proof}

\end{document}
